% Nicolás. S16--S19. Resultados, por qué la hipótesis se cumplió solo en
% parte, conclusión y trabajo futuro. Total previsto: 17:20.
\section{Resultados}
\begin{frame}{Resultados de la tormenta (capacidad 3$\times$)}
  \centering\small
  \begin{tabular}{lcccc}
    \toprule
    Diseño & \shortstack{\emph{Full resync}\\meta $<1\%$} & \shortstack{\emph{Time to recover}\\meta $\leq 30$\,s} & \shortstack{\emph{Staleness} p99\\meta $<1$\,s} & \shortstack{Carga del\\origen} \\
    \midrule
    Binance, snapshot 1$\times$  & 100\% \xmark & 0\,s \okmark & 3.6\,s \xmark & 120\,MB/s \\
    Binance, snapshot 10$\times$ & 100\% \xmark & 38\,s \xmark & 291\,s \xmark & 131\,MB/s \\
    Binance, snapshot 50$\times$ & 100\% \xmark & nunca \xmark & 295\,s \xmark & 316\,MB/s \\
    Estilo Ably                  & 1.9\% \xmark & 78\,s \xmark & 120\,s \xmark & 0.3\,MB/s \\
    \textbf{RelayBook}           & 100\% \xmark & \textbf{0\,s} \okmark & \textbf{1.1\,s} $\approx$ & \textbf{0.3\,MB/s} \\
    \bottomrule
  \end{tabular}\par
  \medskip
  {\small Reconexiones durante el pico: Binance 164\,684 $\cdot$ Ably 215\,531 $\cdot$ \textbf{RelayBook 629}\par}
  \takeaway{RelayBook: 0\,s para recuperarse y 400 veces menos carga en el origen que Binance.}
  \note{Nicolás. 13:25--14:15 (50 s). Esta tabla es con una capacidad de tres veces
  la carga normal. Binance con snapshot barato se recupera al instante, pero
  durante el pico su atraso llega a 3.6 segundos; con snapshot caro tarda 38
  segundos o no se recupera. RelayBook se recupera en cero segundos, con 1.1
  segundos de atraso, y el origen carga 0.3 MB por segundo contra 120 de
  Binance: 400 veces menos. En la línea de abajo, durante el pico Binance tuvo
  164 mil reconexiones y RelayBook 629, porque a los lentos no se les corta.
  El estilo Ably llegó a 120 segundos de atraso, así que esa predicción se
  cumplió. Pero la columna de full resync dice 100 \% para RelayBook. Vamos a
  ver las dos partes de la hipótesis.}
\end{frame}

\begin{frame}{Parte 1: Binance solo cae si el snapshot es caro}
  \centering% S17. Las 18 corridas del contrato de Binance: tiempo hasta que el 99% de los
% clientes vuelve a tener un libro válido tras el pico (observación: 10 min).
\begin{tikzpicture}[x=1cm,y=1cm]
  \path[use as bounding box] (0,0.05) rectangle (13.5,4.6);
  \node[anchor=west,font=\small\bfseries] at (0,4.4) {Tiempo de recuperación de Binance después del pico};
  % Encabezados
  \node[font=\scriptsize\bfseries,anchor=east] at (2.0,3.85) {capacidad $\to$};
  \foreach \c/\h in {0/1.5,1/2,2/3,3/5,4/10,5/15}
    \node[font=\scriptsize\bfseries] at (2.75+\c*1.2,3.85) {\h$\times$};
  \foreach \r/\k in {0/1,1/10,2/50}
    \node[font=\scriptsize,anchor=east,align=right] at (2.0,3.1-\r*0.8) {snapshot \textbf{\k$\times$}};
  % Celdas: columna / fila / valor / relleno / texto
  \foreach \c/\r/\v/\f/\t in {
      0/0/{11\,s}/colReq!25/colReq,   1/0/{0\,s}/colReq!25/colReq,   2/0/{0\,s}/colReq!25/colReq,
      3/0/{10\,s}/colReq!25/colReq,   4/0/{0\,s}/colReq!25/colReq,   5/0/{0\,s}/colReq!25/colReq,
      0/1/nunca/colWake!75/white,     1/1/{303\,s}/colWake!18/colWake, 2/1/{38\,s}/colWake!18/colWake,
      3/1/{0\,s}/colReq!25/colReq,    4/1/{11\,s}/colReq!25/colReq,  5/1/{10\,s}/colReq!25/colReq,
      0/2/nunca/colWake!75/white,     1/2/nunca/colWake!75/white,    2/2/nunca/colWake!75/white,
      3/2/nunca/colWake!75/white,     4/2/{258\,s}/colWake!18/colWake, 5/2/{8\,s}/colReq!25/colReq} {
    \pgfmathsetmacro{\xl}{2.15+\c*1.2}
    \pgfmathsetmacro{\yb}{2.7-\r*0.8}
    \path[fill=\f,draw=white,line width=1.5pt] (\xl,\yb) rectangle ++(1.2,0.8);
    \node[font=\scriptsize\bfseries,text=\t] at (\xl+0.6,\yb+0.4) {\v};
  }
  % Leyenda
  \node[anchor=west,font=\scriptsize] at (2.15,0.55) {%
    \tikz\fill[colReq!25] (0,0) rectangle (0.35,0.2); $\leq$ 30\,s \quad
    \tikz\fill[colWake!18] (0,0) rectangle (0.35,0.2); más de 30\,s \quad
    \tikz\fill[colWake!75] (0,0) rectangle (0.35,0.2); no se recupera en 10\,min};
  % Lectura
  \node[anchor=north west,font=\scriptsize,text width=3.7cm,align=left] at (9.6,3.95) {%
    \textbf{\color{colWake}5 de 18} corridas no se recuperan, todas con snapshot caro y poca capacidad sobrante.\par\medskip
    El origen gastó hasta \textbf{146\,GB} en snapshots que los clientes ya habían \textbf{abandonado} al reintentar, efecto que sostiene la falla.\par\medskip
    Con snapshot barato, siempre se recupera en 11\,s o menos.};
\end{tikzpicture}

  \takeaway{Ubicamos cuándo el contrato es vulnerable; no probamos que Binance falle.}
  \note{Nicolás. 14:15--15:15 (60 s). Parte uno: ¿el contrato de Binance sigue
  degradado tras el pico? Esta grilla son sus 18 corridas: capacidad en
  columnas, costo del snapshot en filas, y en cada celda cuánto tardó en
  recuperarse. Con snapshot barato, la fila de arriba, siempre se recupera en
  11 segundos o menos: no hay falla metaestable. Con snapshot caro y poca
  capacidad sobrante, en rojo oscuro, no se recupera en los diez minutos que
  observamos: 5 de 18 corridas. Ahí se ve el efecto sostenedor del que habló
  Pedro: el origen sigue generando snapshots para clientes que ya los
  abandonaron porque reintentaron, y llegó a gastar 146 GB en eso. Como no
  sabemos cuánto le cuesta a Binance un snapshot, esto dice en qué condiciones
  el contrato es vulnerable, no que Binance falle. La parte uno se cumple solo
  en esas condiciones.}
\end{frame}

\begin{frame}{Parte 2: ¿por qué 100\% de \emph{full resync}?}
  \begin{columns}[c]
    \begin{column}{0.52\textwidth}
      \centering% S18. Niveles visibles distintos que cambian según los mensajes perdidos,
% medidos sobre la grabación. A 1x llegan ~10 mensajes/s; a 10x, ~100.
\begin{tikzpicture}
  \begin{axis}[
    width=7.3cm,height=5.0cm,
    xmode=log,xmin=1,xmax=1000,ymin=0,ymax=5000,
    xtick={1,10,100,1000},xticklabels={1,10,100,1000},
    ytick={0,1000,2000,3000,4000,5000},scaled y ticks=false,
    yticklabels={0,{1\,000},{2\,000},{3\,000},{4\,000},{5\,000}},
    xlabel={mensajes perdidos},ylabel={niveles distintos que cambiaron},
    tick label style={font=\scriptsize},label style={font=\scriptsize},
    axis line style={colNeutro},clip=false]
    \addplot[colWake,thick,dashed,domain=1:1000,samples=2] {2000};
    \node[font=\scriptsize,text=colWake,anchor=south west] at (axis cs:1.1,2000) {libro completo: 2\,000 niveles};
    \addplot[colAcento,line width=1.4pt] coordinates {
      (1,220) (2,374) (5,728) (10,1083) (20,1522) (50,2220)
      (100,2770) (200,3271) (500,3897) (1000,4475)};
    \draw[colReq,densely dotted,thick] (axis cs:10,0) -- (axis cs:10,1083);
    \draw[colWake,densely dotted,thick] (axis cs:100,0) -- (axis cs:100,2770);
    \fill[colReq] (axis cs:10,1083) circle (2.2pt);
    \fill[colWake] (axis cs:100,2770) circle (2.2pt);
    \node[font=\scriptsize,text=colReq,align=left,anchor=north west] at (axis cs:10.5,1050) {1\,s normal\\$\approx$1\,100 niveles};
    \node[font=\scriptsize,text=colWake,align=right,anchor=south east] at (axis cs:95,2850) {1\,s en el pico\\$\approx$2\,800 niveles};
  \end{axis}
\end{tikzpicture}

    \end{column}
    \begin{column}{0.46\textwidth}\small
      \textbf{Full resync 100\%} \xmark\ En el pico, 1\,s de corte son unos 100
      mensajes que tocan $\sim$2\,800 niveles: más que el libro. La regla de
      seguridad manda el snapshot.\par\medskip
      \textbf{Pero sale del relay}, no del origen, y solo 629 veces: la
      conflation evitó cortar a los lentos.\par\medskip
      \textbf{Staleness 1.1--1.3\,s} $\approx$ Un cliente de 80\,KB/s no alcanza un
      stream a 10$\times$ ($\approx$570\,KB/s): no se corta, pero cada entrega
      le llega con $\sim$1\,s de atraso.\par\medskip
      \textbf{Time to recover 0\,s} \okmark
    \end{column}
  \end{columns}
  \takeaway{La mejora vino de \emph{dónde} se recupera y de la conflation, no del tamaño.}
  \note{Nicolás. 15:15--16:35 (80 s). Parte dos: ¿RelayBook cumple las tres metas?
  Cumplió el tiempo de recuperación, quedó cerca en el atraso y falló en full
  resync. Este gráfico explica por qué. En el eje horizontal, cuántos mensajes
  se perdió el cliente; en el vertical, cuántos niveles distintos cambiaron.
  La línea roja es el tamaño del libro, 2 000 niveles. A ritmo normal, un
  segundo son unos 10 mensajes que tocan unos 1 100 niveles: menos que el
  libro, y el resume por estado gana. En el pico a diez veces, un segundo son
  unos 100 mensajes que tocan unos 2 800 niveles, más que el libro entero.
  Entonces se activa la regla de seguridad y el relay manda el snapshot: por
  eso el 100 \%. Lo importante es que ese snapshot sale del relay, no del
  origen, y fueron solo 629. El atraso quedó entre 1.1 y 1.3 segundos por los
  clientes lentos. La conclusión honesta: la mejora de RelayBook no vino de
  recuperaciones más chicas, sino de dónde se recupera y de la conflation.}
\end{frame}

\begin{frame}{Conclusión y trabajo futuro}
  \begin{columns}[T]
    \begin{column}{0.53\textwidth}\small
      \textbf{\color{colAcento}Conclusión}
      \begin{itemize}\setlength{\itemsep}{3pt}
        \item El contrato de Binance \textbf{puede} quedar caído tras un pico,
          pero solo con snapshots caros y poca capacidad (5 de 18).
        \item RelayBook se recupera \textbf{de inmediato}, mantiene la carga
          del origen fija (0.3 vs 120\,MB/s) y no corta a los lentos.
        \item Reanudar por estado ahorra en cortes breves ($-63\%$ a 1\,s) y
          frente a reenviar mensajes (3--84$\times$), pero \textbf{no en un pico}.
        \item Hipótesis: tiempo \okmark, atraso $\approx$, full resync \xmark
      \end{itemize}
    \end{column}
    \begin{column}{0.44\textwidth}\small
      \textbf{\color{colAcento}Trabajo futuro}
      \begin{itemize}\setlength{\itemsep}{3pt}
        \item Medir cuánto cuesta \textbf{de verdad} generar un snapshot.
        \item Probar \textbf{varios instrumentos} y picos reales, no solo la
          grabación acelerada.
        \item Desplegar relays en una \textbf{red real}, con caídas de relays.
        \item Limitar cuánto tiempo guarda el relay las lápidas.
        \item Ajustar la conflation para bajar el atraso de los lentos.
      \end{itemize}
    \end{column}
  \end{columns}
  \vspace{4mm}\hfill{\large\color{colAcento}¿Preguntas?}
  \note{Nicolás. 16:35--17:20 (45 s). Para cerrar. En la hipótesis: tiempo de
  recuperación cumplido, atraso cerca y full resync refutado. El contrato de
  Binance solo queda caído con snapshots caros y poca capacidad, y RelayBook
  se recupera de inmediato sin cargar al origen. Como trabajo futuro: medir
  cuánto cuesta de verdad generar un snapshot, que es la gran incógnita;
  probar varios instrumentos y picos reales; desplegar relays en una red real;
  limitar cuánto guarda el relay las lápidas; y ajustar la conflation para
  bajar el atraso. Muchas gracias. ¿Preguntas?}
\end{frame}
